\section*{Bab \ref{ch:b2:ligand-field} --- Teori Medan Ligan dan Warna}

\begin{solution}{exo:b2:ligand-field:1}
$d^3$: $t_{2g}^3$, CFSE $-1.2\Delta_o$. $d^8$: $t_{2g}^6e_g^2$, CFSE $-2.4 + 1.2 = -1.2\Delta_o$.
\end{solution}

\begin{solution}{exo:b2:ligand-field:2}
\ce{[Fe(H2O)6]^2+}, $d^6$ \omterm{def:b2:ligand-field:spin}{spin tinggi} $t_{2g}^4e_g^2$: empat. \ce{[Fe(CN)6]^4-}, \omterm{def:b2:ligand-field:spin}{spin rendah}
$t_{2g}^6$: tidak ada (\omterm{def:b2:diatomic-mos:magnetism}{diamagnetik}).
\end{solution}

\begin{solution}{exo:b2:ligand-field:3}
600 nm adalah cahaya jingga; kompleks itu tampak biru (biru kehijauan).
\end{solution}

\begin{solution}{exo:b2:ligand-field:4}
$\ce{Cl-} < \ce{H2O} < \ce{NH3} < \ce{CN-}$.
\end{solution}

\begin{solution}{exo:b2:ligand-field:5}
Air: $\Delta_o < P$, \omterm{def:b2:ligand-field:spin}{spin tinggi}, $t_{2g}^3e_g^2$, lima elektron tak berpasangan. Sianida: $\Delta_o > P$,
\omterm{def:b2:ligand-field:spin}{spin rendah}, $t_{2g}^5$, satu elektron tak berpasangan.
\end{solution}

\begin{solution}{exo:b2:ligand-field:6}
Dalam \ce{[CoCl4]^2-} yang tetrahedral, pemisahannya sekitar $4/9$ pemisahan
oktahedral, dan klorida adalah ligan yang lebih lemah daripada air: pita
d--d bergeser ke panjang gelombang yang lebih panjang (jingga ke merah), dan
kompleks itu tampak biru. Kompleks tetrahedral juga menyerap lebih kuat,
karena itu warnanya pekat.
\end{solution}

\begin{solution}{exo:b2:ligand-field:7}
\ce{Ni^2+} adalah $d^8$. Bujur sangkar planar: keempat orbital yang lebih
rendah menampung kedelapan elektron secara berpasangan, dan
$\mathrm d_{x^2-y^2}$ tetap kosong: \omterm{def:b2:diatomic-mos:magnetism}{diamagnetik}. Tetrahedral: $e^4t_2^4$, dua
elektron tak berpasangan dalam himpunan $t_2$.
\end{solution}

\begin{solution}{exo:b2:ligand-field:8}
$\tilde\nu = 1/(500 \times 10^{-7}\,\text{cm}) = \qty{20000}{cm^{-1}}$; $\Delta_o =
0.1196/(500 \times 10^{-9}) = \qty{2.39e5}{J/mol}$, yaitu \qty{239}{kJ/mol}.
\end{solution}

\begin{solution}{exo:b2:ligand-field:9}
\ce{Zn^2+} adalah $d^{10}$ dan \ce{Sc^3+} $d^0$: tidak ada \omterm{def:b2:ligand-field:d-d}{transisi d--d} yang
mungkin (tidak ada tingkat d kosong untuk \ce{Zn^2+}, tidak ada elektron d
untuk \ce{Sc^3+}), sehingga tidak ada serapan di daerah tampak.
\end{solution}

\begin{solution}{exo:b2:ligand-field:10}
Co(III) $d^6$ + enam pasangan elektron bebas amonia: 18 elektron. Dua belas
mengisi keenam \omterm{def:b2:diatomic-mos:bonding}{orbital ikatan}, enam mengisi $t_{2g}$ yang nonikatan (amonia
memberikan medan yang cukup kuat untuk \omterm{def:b2:ligand-field:spin}{spin rendah} dengan kobalt(III)); $e_g^*$
kosong. Semua elektron berpasangan: \omterm{def:b2:diatomic-mos:magnetism}{diamagnetik}.
\end{solution}

\begin{solution}{exo:b2:ligand-field:11}
\ce{CO} adalah donor $\sigma$ yang baik dan, terutama, akseptor $\pi$:
orbital-orbital $\pi^*$-nya yang kosong menurunkan tingkat $t_{2g}$ melalui
\omterm{def:b2:ligand-field:pi-ligands}{donasi balik}, sehingga memperlebar $\Delta_o$. \ce{F-} adalah donor $\pi$, yang
menaikkan $t_{2g}$ dan mempersempit $\Delta_o$. Model muatan titik hanya
melihat muatan dan mendapatkan urutan yang salah.
\end{solution}

\begin{solution}{exo:b2:ligand-field:12}
Tanpa medan, entalpi hidrasi akan mengikuti garis yang mulus (ion-ion
mengecil sepanjang deret). Air memberikan kompleks akua berspin tinggi yang
CFSE-nya nol untuk $d^0$, $d^5$, $d^{10}$
dan terbesar untuk $d^3$ dan $d^8$: stabilisasi tambahan itu menambahkan
dua punuk, yang puncaknya di dekat \ce{V^2+} dan \ce{Ni^2+}, di atas garis
yang melalui \ce{Ca^2+}, \ce{Mn^2+}, \ce{Zn^2+}.
\end{solution}

\begin{solution}{pb:b2:ligand-field:1}
\textbf{1.} $d^3$.
\textbf{2.} $t_{2g}^3$, satu elektron dalam masing-masing ketiga orbital, spin sejajar.
\textbf{3.} Tidak: tiga elektron muat dalam $t_{2g}$ tanpa berpasangan.
\textbf{4.} $-1.2\Delta_o$.
\textbf{5.} Tiga.
\textbf{6.} CFSE-nya besar dan orbital-orbital $e_g^*$, yang mengarah ke
ligan, kosong: melepaskan atau menambahkan ligan menghabiskan sebagian besar
stabilisasi itu.
\textbf{7.} Enam ion oksida pada sudut-sudut oktahedron yang sedikit terdistorsi.
\textbf{8.} $1/(556 \times 10^{-7}) = \qty{17990}{cm^{-1}}$, sekitar \qty{1.80e4}{cm^{-1}}.
\textbf{9.} $0.1196/(556 \times 10^{-9}) = \qty{2.15e5}{J/mol}$: \qty{215}{kJ/mol}.
\textbf{10.} Kuning-hijau (556 nm) dan ungu (405 nm).
\textbf{11.} Merah, dengan sedikit biru: merah tua rubi.
\textbf{12.} Ion $d^3$ mempunyai lebih dari satu konfigurasi tereksitasi
dengan satu elektron dalam $e_g$; tolakan antarelektron memberikan energi
yang berbeda-beda kepadanya, sehingga muncul beberapa pita (dibahas dalam
jilid Tahun 3).
\textbf{13.} $0.1196/(620 \times 10^{-9}) = \qty{1.93e5}{J/mol}$, \qty{193}{kJ/mol}
(\qty{16130}{cm^{-1}}).
\textbf{14.} Rubi: pemisahannya lebih besar.
\textbf{15.} Ke jingga-merah: kini merah yang diserap, dan hijau (dengan
sedikit biru) yang diteruskan.
\textbf{16.} Dalam beril, situs-situs kromium sedikit lebih besar dan
ion-ion oksida lebih jauh, dan lingkungannya berbeda: medannya lebih lemah.
\textbf{17.} Tidak: ion $d^3$ tetap mempunyai tiga elektron tak berpasangan
dalam medan oktahedral mana pun.
\textbf{18.} $d^9$.
\textbf{19.} Kompleks akua menyerap di daerah merah dan jingga melalui
\omterm{def:b2:ligand-field:d-d}{transisi d--d}: cahaya yang diteruskan berwarna biru.
\textbf{20.} Serbuk itu menjadi putih: tanpa ligan air di sekeliling tembaga
(sulfat mengambil tempatnya), medannya lebih lemah dan pitanya bergeser ke
daerah inframerah.
\textbf{21.} Air terikat kembali pada tembaga dan kompleks akua yang biru
terbentuk.
\textbf{22.} Pada panjang gelombang yang lebih pendek: amonia, yang berada
di atas air dalam deret itu, memberikan pemisahan yang lebih besar (biru
keunguan pekat kompleks amin).
\textbf{23.} Tidak: $t_{2g}^6e_g^3$ adalah satu-satunya konfigurasi.
\textbf{24.} Ya: ada satu elektron tak berpasangan.
\textbf{25.} $\boldsymbol{\Delta_o \approx \qty{1.8e4}{cm^{-1}}}$, sekitar
$\boldsymbol{\qty{215}{kJ/mol}}$.
\end{solution}
